% Generated draft: requires mathematical and source review.
\documentclass[11pt]{article}
\usepackage[margin=25mm]{geometry}
\usepackage{fontspec}
\setmainfont{Latin Modern Roman}
\usepackage{amsmath,amssymb,mathtools,unicode-math}
\setmathfont{Latin Modern Math}
\usepackage{enumitem,needspace,xurl,hyperref}
\setlength{\parindent}{0pt}
\setlength{\parskip}{5pt}
\setlength{\emergencystretch}{3em}
\newcommand{\N}{\mathbb{N}}
\newcommand{\Z}{\mathbb{Z}}
\newcommand{\Q}{\mathbb{Q}}
\newcommand{\R}{\mathbb{R}}
\newcommand{\C}{\mathbb{C}}
\newcommand{\F}{\mathbb{F}}
\newcommand{\A}{\mathbb{A}}
\newcommand{\PP}{\mathbb{P}}
\newcommand{\E}{\mathbb{E}}
\newcommand{\eps}{\varepsilon}
\newcommand{\dd}{\,\mathrm{d}}
\begin{document}
\section*{UnsolvedMath ID: 30002181}
\textbf{Source problem code:} catalogue entry
\subsection{Definitions and mathematical statement}
A line bundle \(L\) on a projective variety is \emph{nef} if \(\deg(L|_C)\geq 0\) for every integral projective curve \(C\subseteq X\). It is \emph{ample} if some positive tensor power is very ample. Over an algebraically closed field, a \emph{rational curve} means an integral curve whose normalization is isomorphic to \(\PP^1\).

\textbf{Question as supplied.} For every Calabi--Yau threefold \(X\), does the existence of a nonzero nef line bundle \(L\) on \(X\) which is not ample imply that \(X\) contains a rational curve?

\textbf{Unresolved conventions.} The record does not specify the base field, projectivity, smoothness, or the precise definition of a Calabi--Yau threefold. The definitions above explain the usual projective, algebraically closed setting without asserting that this setting is specified by the source. Moreover, ``nonzero'' could mean \(L\not\cong\mathcal O_X\), or that the numerical class of \(L\) is nonzero. These conditions must not be identified without additional hypotheses. Thus the question retains the supplied terminology and requires these conventions to be resolved before it becomes a fully specified mathematical statement.
\subsection{Short English statement}
Must a Calabi--Yau threefold carrying a nonzero nef line bundle that is not ample contain a rational curve? The base field, Calabi--Yau convention, and meaning of ``nonzero'' require clarification.
\subsection{Sources}
\begin{itemize}
\item Local registry: \texttt{data/unsolvedmath/problems.json}, permanent UnsolvedMath ID \(30002181\), source problem code \texttt{catalogue entry}, in the Oberwolfach Reports Open Problems collection. The supplied record gives proposed year \(2012\) and no named proposer.
\item Supplied report citation: \emph{Komplexe Analysis (2013)}, Oberwolfach Reports, DOI \texttt{10.4171/owr/2012/43}, pp.\ 2597--2656. \url{https://doi.org/10.4171/owr/2012/43}. The supplied background describes these reports as records of mathematical workshops and their open-problem discussions. The citation's \(2013\) designation and its \(2012\) DOI component are retained as supplied.
\item Supplied literature reference: Diverio--Ferretti, ``Oguiso conjecture.'' \url{https://ems.press/content/serial-article-files/43344}.
\item Supplied literature reference: Lazic--Oguiso--Peternell, ``nef line bundles.'' \url{https://arxiv.org/abs/1601.01273}.
\item Supplied literature reference: ``Rational curves and strictly nef divisors.'' \url{https://arxiv.org/abs/2010.12233}.
\item Literature assessment dated exactly \texttt{2026-08-21}: the supplied curation labels the problem partially solved and states that Oguiso's implication is proved when the Calabi--Yau threefold has second Betti number greater than four, with further special-case results, while the unrestricted statement remains unresolved. This is supplied dated curation derived from the OWR record, not an independently verified literature assessment; the hypotheses and scope of the cited results require independent verification.
\item The registry labels its statement \texttt{corrected\_verified} and reports a manual semantic, typographic, and LaTeX review of the cleanup. These are supplied editorial qualifications, not independent verification of the mathematical formulation or current status.
\end{itemize}
\end{document}
